\chapter{Final System Selection and Conclusion}
\label{chap:conclusion}

Based on the extensive validation and calibration phases conducted in the previous chapters, we select the \textbf{Score-Level Fusion of the RBF Support Vector Machine and the Tied Gaussian Mixture Model} as our final, delivered system. This chapter summarizes the evaluation of this chosen model on the unseen Evaluation dataset and formally addresses the concluding laboratory objectives.

\section{Evaluation of the Delivered System}

\subsection*{Question: Compute minimum and actual DCF, and Bayes error plots for the delivered system. What do you observe? Are scores well calibrated for the target application? And for other possible applications?}

Evaluating our final delivered system (the Fused SVM + GMM) on the unseen Evaluation dataset yields exceptional performance. Across our target applications, the system achieved:
\begin{itemize}
    \item At $\tilde{\pi}=0.1$: Minimum DCF = \textbf{0.406}, Actual DCF = \textbf{0.450}
    \item At $\tilde{\pi}=0.5$: Minimum DCF = \textbf{0.166}, Actual DCF = \textbf{0.168}
    \item At $\tilde{\pi}=0.9$: Minimum DCF = \textbf{0.363}, Actual DCF = \textbf{0.368}
\end{itemize}

Observing these metrics and the corresponding Bayes Error plots (see Figure \ref{fig:fusion_vs_gmm} from Chapter 3), it is evident that the scores are highly well-calibrated for the target applications. The Actual DCF closely tracks the Minimum DCF, particularly at the $\tilde{\pi}=0.5$ and $\tilde{\pi}=0.9$ operational points, where the calibration loss is nearly negligible. 

Furthermore, the scores remain well-calibrated for other possible applications. The Actual DCF curve (red) in the Bayes Error plot remains flat and tightly bounds the Minimum DCF curve (blue) across the entire considered log-odds spectrum from $-4$ to $+4$. This proves that our prior-weighted Logistic Regression fusion strategy successfully projected the raw scores into a globally stable, mathematically sound Log-Likelihood Ratio.

\section{Comparison of Actual DCF and Fusion Effectiveness}

\subsection*{Question: Consider the three best performing systems you selected earlier, and their fusion. Evaluate the corresponding actual DCF, and compare the actual DCF error plots. What do you observe? Was your final model choice effective? Would another model / fusion of models have been more effective?}

Comparing the actual DCF profiles of the calibrated Logistic Regression, SVM, GMM, and the Fused system on the evaluation data, we observe that the Fused system consistently achieves the lowest actual operating cost. For example, at $\tilde{\pi}=0.9$, the fused system reaches an actual DCF of $0.368$, outperforming the calibrated GMM ($0.371$), the calibrated SVM ($0.437$), and the Logistic Regression ($0.462$). 

Our final model choice was highly effective. As observed in the comparative Bayes Error plots, the red actual DCF curve of the fused system consistently sits below the curves of the individual constituent models. 

It is highly unlikely that another fusion combination would have been more effective. The success of our system relies on fusing geometrically complementary models. The RBF SVM provides highly accurate, localized, non-linear margin boundaries, while the Tied GMM provides global, probabilistic density contours. Attempting to fuse a strictly linear, under-performing model like the baseline Logistic Regression into this stack would likely introduce estimation noise without adding any novel discriminative geometric information, ultimately degrading the actual DCF.

\section{Calibration Strategy Effectiveness}

\subsection*{Question: Consider again the three best systems. Evaluate minimum and actual DCF for the target application, and analyze the corresponding Bayes error plots. What do you observe? Was the calibration strategy effective for the different approaches?}

When analyzing the three base systems independently, we observe that the calibration strategy (applying an affine transformation learned via 5-fold cross-validation) was unequivocally effective, though the magnitude of the improvement varied depending on the model's native architecture:

\begin{enumerate}
    \item \textbf{Support Vector Machine:} The strategy was overwhelmingly effective. Because SVMs naturally output unbounded geometric margin distances rather than probabilities, the raw scores were catastrophically miscalibrated at extreme priors (e.g., raw actual DCF of $0.631$ at $\tilde{\pi}=0.1$). Calibration successfully mapped these distances into true LLRs, dropping the actual DCF to $0.347$ on the validation set and stabilizing the Bayes Error plot.
    \item \textbf{Logistic Regression:} The strategy was highly effective at correcting the bias introduced by the model's initial training parameters. The calibration successfully shifted the threshold to align with the target applications, closing the gap between the actual and minimum DCF curves at the extremes.
    \item \textbf{Gaussian Mixture Model:} The strategy was effective, though less drastic. Because the GMM natively computes probabilistic likelihoods, its baseline calibration was already strong. However, the calibrator successfully compensated for minor density distortions caused by the non-Gaussian skewed tails in the dataset, further refining the actual DCF across all target applications.
\end{enumerate}

In conclusion, score calibration and geometrically diverse score-level fusion are indispensable steps for deploying generative and discriminative pattern recognition systems in real-world, highly asymmetric target applications.